% chapters/appendix_c_reward.tex
\chapter{Reward Function Specification}
\label{app:reward}

The immediate move rewards and event-driven rewards are specified in \cref{tab:sys:move_rewards,tab:sys:event_rewards}. A linear distance penalty $r_{\text{dist}} = -0.001 \times d_{\text{metres}} - 0.0005 \times t_{\text{seconds}}$ is applied to every executed move, where $d_{\text{metres}}$ is the Manhattan distance the crane travels and $t_{\text{seconds}}$ is the estimated travel time. This penalty is small relative to the type bonuses but accumulates over many moves, providing a gradient signal toward compact container placement.

\section{Reward Design Rationale}
\label{app:reward:rationale}

The reward function balances three competing objectives:

\begin{enumerate}
  \item \textbf{Throughput:} High rewards for productive moves and train completions incentivise maximising container flow through the terminal.
  \item \textbf{Service quality:} Truck fast-service bonuses reward timely dispatch without penalising the agent for factors outside its control.
  \item \textbf{Yard order:} Restack penalties and EOD blocking penalties discourage short-sighted moves that create future problems.
\end{enumerate}

The magnitudes were calibrated empirically during development. The train departure bonus ($+10.0$) is the largest single reward, reflecting the irrecoverable nature of missed train loads. The IDLE penalty ($-0.1$) is large enough to prevent the agent from learning an ``always idle'' policy but small enough that the agent can learn to wait when no productive action is available. The restack penalty ($-0.2$) is deliberately mild: some restacking is operationally necessary, and a large penalty would prevent the agent from learning multi-step restack-then-load sequences.

\subsection{Reward Shaping Corrections}
\label{app:reward:corrections}

During Phase~2 evaluation, two reward shaping defects were identified that systematically punished the agent for correct behaviour, preventing effective learning at scale:

\begin{enumerate}
  \item \textbf{End-of-day leftover penalty (all containers $\to$ missed only).}
    The original implementation penalised \emph{every} container remaining in the yard at end of day ($-0.3$ each), including legitimate carryover inventory awaiting future pickup.
    At 100+ containers per day, the yard naturally accumulates 100--200 containers whose departure dates lie days or weeks in the future; penalising all of them produced $-30$ to $-60$ of unearned negative reward per day.
    The corrected penalty counts only containers whose \texttt{departure\_date} is on or before the current day---i.e.\ containers that \emph{should have left} but did not.

  \item \textbf{Continuous truck waiting penalty (removed).}
    A per-truck-per-minute penalty ($-0.01 \times n_\text{trucks} \times \Delta t_\text{min}$) was applied on every simulation step where all cranes were busy.
    Because truck count ($n_\text{trucks}$) includes already-parked and actively-served trucks---not only those genuinely waiting---this penalty accumulated to $-100$ to $-200$ per day at 100 containers, far exceeding the positive reward from productive moves.
    The penalty was removed; the truck departure reward ($+2.0$ for fast service, $+0.0$ for slow) provides a sufficient outcome-based incentive without punishing the agent for throughput limitations of the crane system.
\end{enumerate}

Together, these two corrections shifted the expected daily reward from deeply negative (even for competent policies) to a range where positive reward is achievable, enabling meaningful gradient signal during Phase~2 training.

\subsection{Missed Container Removal}
\label{app:reward:missed-removal}

Containers whose \texttt{departure\_date} has passed are removed from the yard at the end of each simulated day, after the EOD penalty has been applied.
This design decision is motivated by three considerations:

\begin{enumerate}
  \item \textbf{Operational realism.}
    In a real intermodal terminal, containers that miss their scheduled departure are not left to occupy prime stacking positions indefinitely.
    The terminal's dispatch team reschedules them on the next available train or truck, relocates them to a holding area, or flags them for customer intervention.
    The simulation abstracts this exception handling process as removal at end of day.

  \item \textbf{Agent scope.}
    The agent controls crane movements within a single day.
    It is already penalised for a missed departure at the moment the train leaves ($-3.0$ per container, \cref{tab:sys:event_rewards}).
    Retaining the missed container in the yard penalises the agent \emph{again} on every subsequent day for a failure it can no longer correct.
    An agent should only receive penalties for outcomes it can still influence; compounding irrecoverable penalties distorts the learning signal without providing actionable information.

  \item \textbf{Markov property.}
    DQN assumes that the current state contains sufficient information for optimal future decisions.
    Accumulated dead inventory violates this assumption: the agent cannot distinguish a container awaiting a future train from one that missed its train three days ago and will never leave.
    Removing missed containers ensures that every container in the state tensor is one the agent can meaningfully act on, preserving the representational integrity of the Markov state.
\end{enumerate}

Without removal, missed containers compound across days: they consume yard capacity (causing destination masks to become empty and move attempts to fail), accumulate EOD penalties on every subsequent day, and create stacking inversions with newly arrived containers.
In testing, this compounding effect caused daily rewards to drop from $+35$ on Day~1 to $-1{,}479$ by Day~4 despite improving operational metrics, because the yard filled with irrecoverable dead inventory.
